\changecolours{ScoutGreen}
\section[Referências]{Referências Bibliográficas}

% ------------------------------------------------------------------------------
\begin{frame}{Referências Bibliográficas Fundamentais}
  \alert{Obras Canônicas do Modelo Relacional e Otimização}\par
  \vspace{0.3cm}
  \scriptsize
  \begin{enumerate}
    \item \textbf{CODD, Edgar Frank.} \emph{A Relational Model of Data for Large Shared Data Banks}. Communications of the ACM (CACM), New York, v. 13, n. 6, p. 377–387, June 1970.
    \item \textbf{CODD, Edgar Frank.} \emph{Relational Completeness of Data Base Sublanguages}. In: RUSTIN, R. (Ed.). \emph{Data Base Systems}. Englewood Cliffs: Prentice-Hall, 1972. p. 65–98.
    \item \textbf{ELMASRI, Ramez; NAVATHE, Shamkant B.} \emph{Sistemas de Banco de Dados}. 7. ed. São Paulo: Pearson Education do Brasil, 2018. (Capítulo 8: Álgebra Relacional e Cálculo Relacional; Capítulo 19: Algoritmos para Otimização de Consultas).
    \item \textbf{HEUSER, Carlos Alberto.} \emph{Projeto de Banco de Dados}. 6. ed. Porto Alegre: Bookman, 2009. (Capítulo 5: Linguagens Formais de Consulta e Otimização).
    \item \textbf{SILBERSCHATZ, Abraham; KORTH, Henry F.; SUDARSHAN, S.} \emph{Sistema de Banco de Dados}. 7. ed. Rio de Janeiro: GEN/LTC, 2020. (Capítulo 6: Linguagens Formais; Capítulo 13: Otimização de Consultas).
    \item \textbf{DATE, Chris J.} \emph{Introdução a Sistemas de Bancos de Dados}. 8. ed. Rio de Janeiro: Elsevier Campus, 2004. (Capítulo 7: Cálculo Relacional; Capítulo 18: Otimização).
  \end{enumerate}
\end{frame}
